\section{Results}
\label{sec:results}

\subsection{Overview of INTENSE}
\label{sec:overview_intense}

The INTENSE framework
detects which neurons encode specific behavioral or environmental variables from
large-scale calcium imaging data (\fig{fig:method}). INTENSE works directly with
raw calcium fluorescence, eliminating the need for spike deconvolution. It uses
mutual information --- which captures any statistical dependency, linear or
nonlinear, equaling zero if and only if variables are independent
\cite{Shannon1948, Borst1999, Strong1998} --- as its core measure. Specifically,
INTENSE employs the Gaussian Copula MI (GCMI) estimator
\cite{gcmi, ma2011mutual}, which provides efficient closed-form MI computation
(see Methods: Section~\ref{sec:mi_computation}), making it feasible to analyze
thousands of neuron-behavior pairs across large populations.

Statistical significance is assessed
through circular-shift permutation testing with parametric p-value modeling
(\fig{fig:method}E; see Methods: Section~\ref{sec:statistical_assessment}).
Systematic delay optimization accounts for calcium indicator kinetics and
prospective or retrospective encoding
(Section~\ref{sec:temporal_alignment}).

\begin{figure}
  \centering
  \includegraphics[width=1\linewidth]{figs/fig1_delegated.png}
  \caption{\textbf{A}: Examples of rescaled $dF/F$ calcium fluorescence traces (blue) with detected events (grey).
\textbf{B}: Examples of behavioral/environmental
variables used for correlation with neuronal activity: continuous (head direction, top),
discrete (object interaction, middle), multidimensional (animal coordinates,
bottom). \textbf{C}: Schematic overview of INTENSE pipeline for quantifying neuronal selectivity. \textbf{D}: Left: example object-interaction
cell: fluorescence trace (blue) and interaction periods (green) . Right: distributions of scaled $dF/F$ values inside (green) vs
outside (grey) object interaction periods. \textbf{E}: Distribution of
neuron-behavior pairs in the ``significance-power'' space with
marginal distributions. Pairs considered relevant are shown in green.}

  \label{fig:method}
\end{figure}

Beyond detection, INTENSE provides tools for analyzing mixed selectivity
(\fig{fig:results}F,G). Many hippocampal and cortical neurons show apparent
selectivity to multiple behavioral variables \cite{mix2,mix3}, but distinguishing
true mixed selectivity from correlations induced by covarying behaviors is essential
for understanding neural coding principles and the true dimensionality of
representations. The framework's disentanglement procedure uses conditional MI and interaction
information \cite{ghassami2017} to separate genuine multi-variable encoding from
spurious associations (Section~\ref{sec:disentangling_selectivity},
Section~\ref{sec:intense_deciphers}).

We validated INTENSE on synthetic datasets with known ground truth across a wide
range of signal-to-noise ratios and response reliability conditions
(\fig{fig:synthetics}, Section~\ref{sec:synthetic_validation}), and against
classic event-based place cell detection methods
(\fig{fig:pc}, Section~\ref{sec:intense_validates}), then applied it to
hippocampal calcium imaging data from freely behaving mice, recovering diverse
functional cell types including place cells, head direction cells, and neurons with
genuine mixed selectivity (\fig{fig:results}A,B,D,E,
Section~\ref{sec:intense_deciphers}).


\subsection{INTENSE validates against established place cell detection methods}
\label{sec:intense_validates}

\begin{figure}
  \centering
  \includegraphics[width=0.9\linewidth]{figs/fig2_delegated.png}
  \caption{\textbf{A}: Upper -- Continuous Wavelet Transform map (CWT map) of a
fluorescence signal. Cross-scale ridges correspond to detected events. Lower -- fluorescent trace (blue) and detected events (grey). \textbf{B}: Example animal
trajectory, colors represent rescaled calcium activity of a single place
cell. Stars represent detected events. \textbf{C}: Left -- intersection of
identified PC populations from INTENSE (red) and classic analysis (blue) with
$p<0.05$. Right: the same for all cells identified as PC by any method. \textbf{D}: Example of activity maps computed from
raw calcium signals (left) and from detected events (right) with high SSIM score.
\textbf{E}: Similarity metrics between activity maps computed from raw calcium
signals and from detected events as functions of confidence (measured via
$p$-value). Left: mean squared error, center: peak signal-to-noise ratio, right:
structural similarity score. Shaded regions indicate confidence intervals.
Vertical lines show the $p=0.001$ threshold, as in \fig{fig:pc}F. \textbf{F}: Relative intersection between PC populations computed
via INTENSE and classic analysis as a function of confidence (measured via
$p$-value). Intersection is computed on cells for which $p_{INTENSE} \leq p$ and
$p_{classic} \leq p$ simultaneously. Neurons considered as ``non-PC'' by both
analyses are excluded. Shown are intersections as fractions of respective
populations (red for INTENSE, blue for classic); Jaccard coefficient
$J$=(INTENSE $\bigcap$ Classic) / (INTENSE $\bigcup$ Classic), orange; proportion
of all identified PCs (by any method) left after thresholding by $p$.}
  \label{fig:pc}
\end{figure}


To verify the precision of INTENSE, we applied it to the well-established task of
identifying place cells (PCs)---neurons that exhibit selectivity for specific locations in
the environment \cite{o1971hippocampus}.

PC detection provides a challenging validation benchmark, as the choice of
classification method greatly determines the population of cells identified
\cite{Grijseels2021}, with overlap as low as 30--40\% even for different
event-based methods.

We compared INTENSE, applied directly to calcium fluorescence signals and animal
coordinates (see Methods: Section~\ref{sec:intense_pc_identification}), with a
classic event-based method that constructs binarized activity maps from deconvolved
calcium events \cite{skaggs1993information} (see Methods:
Section~\ref{sec:classic_pc_identification}). To enable a direct, like-for-like comparison, we expressed both methods’ outputs as (i) a binary place-cell label and (ii) an associated confidence score (p-value) derived from each method’s native null model (see
Methods: Section~\ref{sec:comparison_pc_populations}).

Despite modest overlap ($\sim$40\%) when considering all cells identified by either
method, the agreement increased dramatically when confidence thresholds were applied
(\fig{fig:pc}C). The overlap grew monotonically with threshold stringency, reaching
90\% at $p < 0.001$---a threshold that retained 95\% of the total place cell
population (\fig{fig:pc}F). Further threshold increases yielded near-perfect
agreement, though at the cost of excluding many putative place cells.

Significant heterogeneity of PC firing rates and stability has been established both
experimentally \cite{Mizuseki2012} and theoretically \cite{Lian2022}.
Our results indicate that INTENSE and spike-based methods primarily differ in their
treatment of ``weak'' place cells, but converge on the same ``core'' population when
statistical confidence requirements are imposed. Thus, our general
information-theoretic approach successfully captures the spatial coding properties
identified by specialized algorithms.


To verify that all matching PCs identified by both methods have the same pattern of
spatial selectivity, we analyzed the coincidence of activity maps. There exists a
wide variety of place field detection methods \cite{Grieves2017}, but they are
highly dependent on the data type. To overcome this difficulty, we compared
discretized activity maps (normalized by time spent in the bins; see Methods:
Section~\ref{sec:comparison_spatial_selectivity} and \fig{fig:pc}D).

We quantified the correspondence between activity maps using structural
similarity (SSIM) and complementary metrics (see Methods:
Section~\ref{sec:comparison_spatial_selectivity}).

All spatial correspondence metrics improved with increasing confidence thresholds
(\fig{fig:pc}E), confirming that place cells identified with high confidence by both
methods exhibit highly similar spatial firing patterns. This multi-metric
validation demonstrates that INTENSE not only identifies overlapping cell
populations but also accurately captures their functional properties.


\subsection{INTENSE reveals diverse neuronal selectivity}
\label{sec:intense_deciphers}

\begin{figure}
  \centering
  \includegraphics[width=1\linewidth]{figs/fig3_delegated.png}
  \caption{\textbf{A}: Example ``object interaction neuron'' and $dF/F$ distributions during interaction (green) and
control periods (grey). \textbf{B}:
Example ``rear neuron'' and $dF/F$ distributions during rearing (orange) and control periods (grey). \textbf{C}: Left -- distribution of neurons selective to 1–3 variables. Right -- distribution of optimal delays: negative ($-$), near-zero and positive
($+$). Statistical comparisons are detailed in
Methods. \textbf{D}: Head direction tuning map of a multi-selective neuron.
\textbf{E}: 2D spatial activity map of this neuron. \textbf{F}: Heatmap of MI between
significantly correlated behavioral features (behavior-behavior analysis). Connection strength is represented by color intensity.
\textbf{G}: Schematic of the interdependencies between information measures in
neurons with entangled selectivity.}

  \label{fig:results}
\end{figure}


Having validated INTENSE against established methods, we applied it to discover
novel neuron-behavior associations in our hippocampal recordings. We analyzed
associations between neural activity and a comprehensive behavioral repertoire
including spatial position, head direction, movement dynamics, and discrete
behavioral states (see Methods: Section~\ref{sec:behavioral_analysis}). Following
identification of significant neuron-behavior pairs, we performed
``disentanglement'' analysis to distinguish true mixed selectivity from spurious
associations arising from behavioral correlations (\fig{fig:results}A--E; see
Methods: Section~\ref{sec:disentangling_selectivity}).

The disentanglement procedure uses conditional mutual information (CMI) to
separate primary from secondary selectivities, and interaction information (II) to
classify feature pairs as redundant or synergistic (see Methods:
Section~\ref{sec:disentangling_selectivity}).

The results showed the presence of multiple neurons that are selective for various
specific behavioral features (such as object interaction, rearing, etc., as
shown in \fig{fig:results}A,B).
% [AUTHOR: The following paragraph on probabilistic activation patterns needs review.
% Consider whether this interpretive discussion belongs in Results or Discussion.]
These behaviorally selective neurons showed
probabilistic activation patterns---they neither fired deterministically during
every instance of their preferred behavior nor remained completely silent otherwise.
This can be explained by the probabilistic nature of neural coding and the need for
more detailed identification of behavioral acts based on the behavioral context.
Acts that may seem identical from the perspective of the researcher based on the
kinematic characteristics of the reference points, may have different internal
meanings for the animal.

Alongside discrete behavior encoding, we found numerous cells selective for
continuous variables: movement speed, head direction, and spatial position
(\fig{fig:results}D,E). Cells selective for locomotion state, head direction or speed have been increasingly recognized in
hippocampus \cite{kropff2015speed,Gois2018,Acharya2016}.

\paragraph{Temporal alignment of neural activity and behavior.}
\label{sec:temporal_results}
To account for temporal offsets between neural activity and calcium
fluorescence, INTENSE optimizes the delay between signals (see Methods:
Section~\ref{sec:temporal_alignment}).
Optimal delays were binned into three categories (negative, near-zero,
and positive) for statistical analysis. The majority of significant associations
occurred at positive delays (44.5\%), consistent with the expected temporal lag of
calcium indicators, with 41.0\% at near-zero delays and only 14.5\% at negative
delays (\fig{fig:results}C). Because our primary goal was to characterize associations consistent with calcium-indicator lag, we report delay distributions for all significant pairs and, unless stated otherwise, restrict downstream selectivity analyses to pairs with non-negative optimal delays. 

\paragraph{Quantifying and disentangling mixed selectivity.}
Of the 30{,}105 neuron-sessions analyzed across 16 mice and 4 sessions, 6{,}217
(20.7\%) showed significant selectivity to at least one behavioral variable.
Among these, 1{,}377 (22.1\% of selective neurons) responded to two or more
variables before disentanglement (\fig{fig:results}C).

Pairwise significance analysis of behavioral variables revealed tightly coupled
clusters (\fig{fig:results}F). Locomotion-related variables (speed, rest,
walk, locomotion state) showed significant co-selectivity in 82--100\% of
mouse-session combinations, indicating strong behavioral correlations that could
inflate apparent mixed selectivity. Similarly, spatial zone variables (corners,
walls, center) formed a distinct cluster (77--91\% co-significance).

The CMI-based disentanglement procedure resolved a substantial fraction of these
multi-variable associations. After controlling for behavioral correlations, the
number of multi-selective neurons decreased from 1{,}377 to 949 (15.7\% of
selective), with approximately a third of multi-variable associations explained by
behavioral redundancy. The effect was most dramatic for speed selectivity: of 201
neurons initially classified as speed-selective, only 7 retained significance after
conditioning on locomotion state---consistent with speed and locomotion sharing
100\% co-selectivity across all sessions.

Interaction information provided a complementary summary of behavioral variable associations.
For the speed--locomotion pair, II was consistently negative (redundant), confirming
that these variables carry overlapping information about neural activity.
Conversely, pairs such as body direction and head direction showed near-zero or
positive II in most sessions, consistent with these closely related variables
carrying largely overlapping information about neural activity.

Even after disentanglement, a substantial number of neurons retained significant
associations with multiple variables. For example, the neuron illustrated in
\fig{fig:results}D,E showed genuine mixed selectivity for speed, head direction, and
position simultaneously. While such multiply-selective neurons comprise a minority
of the population (\fig{fig:results}C), they may serve as crucial integration points
linking different functional systems \cite{mix3}.

\paragraph{Mutual information detects associations invisible to linear correlation.}
\label{sec:mi_vs_corr}
Replacing mutual information with Pearson correlation in the otherwise identical INTENSE pipeline detected 2.2 times fewer selective neurons and 2.9 times fewer significant neuron--feature pairs. The gap was largest for discrete behavioral variables (5--11$\times$) and remained substantial even for continuous variables such as speed (3.2$\times$). These results confirm that a model-free, information-theoretic metric substantially expands the repertoire of detectable neural selectivities.

\subsection{Synthetic validation demonstrates INTENSE robustness}
\label{sec:synthetic_validation}

\begin{figure}
  \centering
  \includegraphics[width=1\linewidth]{figs/fig4_delegated.png}
  \caption{\textbf{A}: Stages of generating pseudo-calcium signal, associated with
a continuous variable. 1 -- continuous behavioral time series with selected range of
interest; 2 -- randomized binary time series (with $p_{skip}$); 3 -- heterogeneous
Poisson process with two different rates; 4 -- resulting pseudo-calcium signal after
convolution with a calcium indicator kernel. For discrete variables, the procedure
starts from stage~2.
\textbf{B}: $F_1$-score heatmaps for the continuous variable test across the whole
parameter range. Methods shown: correlation-based, MI-based, correlation-based with
shuffles, INTENSE. \textbf{C}: Same as B for discrete variables (generated starting
from stage~2 of A). Methods shown: average-based, MI-based,
average-based with shuffles, INTENSE. \textbf{D}: Precision-Recall maps for
continuous variable test. Methods shown: correlation-based, MI-based,
correlation-based with shuffles, INTENSE. Left -- precision-recall coordinates of
all methods for different SNR values (taken at $p_{skip} = 0$). Right --
precision-recall coordinates of all methods for different $p_{skip}$ values (taken
at $SNR = 64$). Vertical lines represent random guesser precision ($PR_{random}
=0.05$). \textbf{E}: Precision-Recall maps for discrete variable test. Methods
shown: average-based, MI-based, average-based with shuffles, INTENSE. Left --
precision-recall coordinates of all methods for different SNR values (taken at
$p_{skip} = 0$). Right -- precision-recall coordinates of all methods for different
$p_{skip}$ values (taken at $SNR = 64$). Vertical lines represent random guesser
precision ($PR_{random} =0.05$).}
  \label{fig:synthetics}

\end{figure}


To rigorously evaluate INTENSE's detection capabilities, we conducted systematic
tests on synthetic datasets with known ground truth. We generated experiments
containing 20 behavioral time series and 500 pseudo-calcium signals, with each
calcium signal linked to exactly one behavioral variable (5\% prevalence of true
associations).

\subsubsection{Synthetic framework and the role of temporal controls}
\label{sec:bio_synthetic}

We generated realistic synthetic experiments where calcium signals were
linked to behavioral variables through a biologically plausible process
(\fig{fig:synthetics}A; see Methods: Section~\ref{sec:synthetic_data_generation}).
The framework varied signal-to-noise ratio
(SNR $\in \{2, 4, 8, 16, 32, 64\}$) and response reliability
($p_{\text{skip}} \in [0.0, 0.8]$).

A method comparison revealed a fundamental challenge: temporal autocorrelation
creates overwhelming false positive rates without proper statistical controls
(\fig{fig:synthetics}B--E). Methods using standard correlation or mutual information
tests without shuffle controls achieved precision of only 8.7\% and 6.4\%
respectively for continuous variables, barely exceeding the 5\% random baseline.
This demonstrates why shuffle-based significance testing is essential---the slow
dynamics of both calcium signals and behavioral variables create spurious
correlations that appear highly significant under naive statistical tests.

\subsubsection{INTENSE outperforms other shuffle-based approaches}
\label{sec:intense_outperforms}

Among methods employing shuffle controls, INTENSE demonstrated clear advantages over
correlation-based approaches. Although the correlation-based method with shuffle controls achieved high
precision (100.0\% at SNR = 64), it suffered from extremely poor recall
(maximum 4.5\%) across
---detecting less than 5\% of true associations even under optimal
conditions. This pattern reveals a fundamental limitation: correlation-based
methods, even with proper statistical controls, remain restricted to linear
associations and miss the complex nonlinear dependencies often present in neural
data.

In contrast, INTENSE maintained strong performance across the full spectrum of
signal strengths. With perfect response reliability ($p_{\text{skip}} = 0$),
precision increased from 44.2\% at SNR = 2 to 85.5\% at SNR = 64, while recall
improved from 6.8\% to 41.7\% for continuous variables.
Discrete behavioral variables yielded substantially better results: precision rose from 73.5\% to
94.6\%, with recall increasing dramatically from 24.3\% to 100.0\% across the same
range.

\subsubsection{Two-stage testing balances efficiency and power}
\label{sec:two_stage_testing}

The two-stage testing procedure (see Methods: Section~\ref{sec:statistical_assessment})
proved highly effective at managing the
computational burden of large-scale analysis. Stage 1 identified 3--6\%
of all possible pairs as candidates.
Precision improved from 62.0\% to 85.5\% for continuous variables and from 89.4\% to 94.6\%
for discrete variables between stages, while recall decreased only modestly, demonstrating
effective false discovery control.

\subsubsection{Performance under biological variability}
\label{sec:performance_variability}

Real neurons show variable response reliability---they may skip responses due to
adaptation, internal state changes, or stochastic factors. INTENSE demonstrated
graceful performance degradation as response reliability decreased. At SNR = 64,
precision for continuous variables decreased from 85.5\% at $p_{\text{skip}} = 0$ to
74.1\% at $p_{\text{skip}} = 0.5$. Discrete variables showed greater robustness,
maintaining 94.8\% precision with 99.7\% recall even at $p_{\text{skip}} = 0.5$,
demonstrating robustness to moderate response unreliability commonly observed in real neural data.

An intriguing phenomenon emerged at very high signal-to-noise ratios: raw mutual
information and average-based metrics exhibited performance degradation due to calcium
indicator saturation effects (\fig{fig:synthetics}B,C). When neurons fire frequently
and strongly, slow decay dynamics cause fluorescence to remain elevated between events,
merging individual episodes into sustained plateaus. For discrete variables, this caused
raw MI precision to drop from 92.9\% at SNR = 4 to 70.0\% at SNR = 64, and
average-based detection with shuffle controls to collapse from F1 = 51.9\% to just
10.1\% over the same range. INTENSE proved robust to these effects, with F1 improving
monotonically from 95.2\% to 97.2\%. This robustness arises because circular-shift
shuffles construct the null distribution from the same saturated signals, so saturation
elevates both true and null MI values equally, preserving the validity of the
significance test.

These synthetic results demonstrate that INTENSE successfully addresses the key
challenges in detecting neuron-behavior associations: controlling false discoveries
from temporal structure, capturing complex associations beyond linear correlation,
and maintaining sensitivity despite biological variability. The framework's robust
performance across diverse conditions establishes confidence for its application to experimental data where ground truth is unknown. 

